\section{Arrangements and affine oriented matroids}\label{sec:background}

\subsection{Affine and central arrangements}
Write $\C=\{H_i\}_{i\in E_0}$, where
$H_i=\{f_i=0\}$ and $f_i(x)=\ell_i(x)+c_i$ for linear functionals $\ell_i \in (\R^r)^*$ and scalars $c_i \in \R$. To $\C$ we attach the \emph{central arrangement} $\C_0=\{H_i^0\}_{i\in E_0}$ where $H_i^0=\{\ell_i=0\}$. Moreover, by homogenizing the defining equations of $\C$ we obtain the central arrangement $\widehat\C$ in $\R^{r+1}$ with hyperplanes defined by $\widehat f_i(x,t)=\ell_i(x)+c_it$ for $i\in E_0$ and an additional hyperplane at infinity defined by $\widehat f_0(x,t)=t$ and indexed by $0$. The slice $t=1$
recovers $\C$, and restriction of the forms $\widehat f_i$, $i\in E_0$,
to $t=0$ gives $\C_0$.

To a central arrangement, one may associate a matroid defined on the set indexing the hyperplanes whose independent sets correspond to the linearly independent subsets of the linear forms defining the hyperplanes.
This gives the matroids $M_0$ and $M$ associated to $\C_0$ and
$\widehat\C$, on $E_0$ and $E=E_0\sqcup\{0\}$, respectively.
The assumption that $\C$ is essential means that the forms $\ell_i$ span $(\R^r)^*$.
Thus $M_0$ has rank $r$, $M$ has rank $r+1$, and $M_0$ is the contraction of $M$ by $0$.

\subsection{Matroid conventions}
For a matroid $N$, write $\rk_N J$ for the rank of a subset $J$.
A \emph{flat} is a subset whose rank increases whenever an element
outside it is added. Let $N^J=N|J$ and $N_J=N/J$ denote restriction
and contraction, and let $N^*$ be the dual matroid. Write $\cL(N)$
for the lattice of flats and $\Z(N)$ for its cyclic flats, namely
the flats $J$ for which $N^J$ has no coloops. 

An affine matroid $(M,0)$ is a matroid on a ground set $E$ with a
choice of distinguished element $0\in E$. For such a pair, we write
\[
 \cL=\cL(M),\qquad \cL_0=\{F\in\cL\mid 0\in F\}.
\]
The flat correspondence for contraction identifies $\cL(M_0)$ with
$\cL_0$ by $K\mapsto K\cup\{0\}$. We distinguish these flats as
subsets of $E_0=E\setminus\{0\}$ and $E$, respectively.


For a loopless matroid $N$ of rank $s$, let
\[
 \chi_N(t)=\sum_{F\in\cL(N)}\mu(\varnothing,F)t^{s-\rk_N F},
\]
where $\mu$ is the M\"obius function of $\cL(N)$. We use
$\mup(N)=(-1)^s\chi_N(0)$ and, for $s>0$,
$\beta(N)=(-1)^{s-1}\chi_N'(1)$. For a matroid with loops we put
$\chi_N=0$ and $\mup(N)=\beta(N)=0$. For the empty matroid,
$\mup(\varnothing)=1$ and $\beta(\varnothing)=0$.

Returning to the affine matroid $(M,0)$, recall that for a flat $F \in \cL$ we set $w_F=\beta(M^F)\beta(M_F)$.

\subsection{Affine oriented matroids}
Central and affine arrangements are generalized by oriented matroids
and affine oriented matroids. These axiomatize the sign vectors of
faces, similarly to the way matroids axiomatize linear independence
among vectors.
An oriented matroid coming from a real hyperplane arrangement is
called \emph{realizable}. Not all oriented matroids are realizable.
Conjecturally, our determinant theorems extend to affine oriented matroids.
We recall the terminology.

An \emph{oriented matroid} $\M$ on a finite ground set $E$ is specified
by sign vectors $G\in\{0,+,-\}^E$, called \emph{covectors}, satisfying
the covector axioms \cite{OrientedMatroids1999}.
We write $M$ for its underlying matroid and assume that $M$ is simple
of rank $r+1$. An \emph{affine oriented matroid} $(\M,0)$ has a
distinguished element $0\in E$. A covector $G$ is \emph{feasible} if
$G(0)=+$. A \emph{tope} is a covector with no zero entries, and we
denote the feasible topes by $\cF$.

The homogenized arrangement $\widehat\C$ yields an affine oriented matroid $(\mathcal{M},0)$ on the ground set $E=E_0\cup\{0\}$ whose covectors are the sign vectors
\[
 \bigl(\operatorname{sign}\widehat f_i(x,t)\bigr)_{i\in E},
 \qquad (x,t)\in\R^{r+1}.
\]
The feasible covectors (topes) of $(\mathcal{M},0)$ correspond exactly to the sign vectors of all faces (chambers) in the affine arrangement $\C$.
Similarly, we can consider the oriented matroid $\mathcal M_0=\mathcal{M}/0$ which corresponds to the central arrangement $\C_0$.

Many combinatorial aspects of arrangements translate to oriented matroids.
The face order and zero set are given by
\[
 G\leq F\ \Longleftrightarrow\ G(i)\in\{0,F(i)\}\text{ for every }i,
 \qquad z(G)=\{i\mid G(i)=0\}.
\]
For every nonzero covector $G$, define
\[
 \dim G=r-\rk_M z(G).
\]
For feasible $G$, this agrees with the dimension of its affine face
in the realizable case. 

A feasible face $G$ is \emph{bounded} if
every nonzero covector $F\leq G$ is feasible. We denote the bounded
feasible topes by $\cB\subseteq\cF$. These correspond to the bounded
chambers of the affine arrangement.  As recalled in~\cite[Proposition 3.5]{Lam}, in both the case of hyperplane arrangements and more generally for affine oriented matroids, the number of bounded-regions (or bounded feasible topes) is given by
\[
 |\cB|=\beta(M).
\]


For an arrangement, write $H_J=\bigcap_{i\in J}H_i$.
We use this geometric notation for oriented matroids as well. 
 For $J\subseteq E_0$, write
$P^J$ for the sign vector obtained from $P$ by reversing the signs in
$J$, and abbreviate $P^{\{i\}}$ to $P^i$.

We also use the composition of sign vectors
\[
 (F\circ G)(i)=
 \begin{cases}F(i),&F(i)\ne0,\\G(i),&F(i)=0.\end{cases}
\]
Covectors are closed under composition. For an arrangement, this fills
the zero entries of $F$ according to $G$, describing a small displacement
from the face $F$ towards the face $G$.

We also use the \emph{strong elimination axiom} in the following
special case. If two covectors agree outside a coordinate $i$ and
have opposite signs at $i$, then setting that coordinate to zero
and keeping their common entries gives a covector.
For hyperplanes, this follows by intersecting the segment between
points of the two faces with $H_i$.

These definitions give the following dictionary.

\begin{center}
\small
\begin{tabularx}{\textwidth}{@{}XX@{}}
\toprule
Affine hyperplane arrangement & Affine oriented matroid\\
\midrule
Chamber $P\in\cF$ & Tope $P$ with $P(0)=+$\\
Face in $\bar P$ & Covector $G\leq P$ with $G(0)=+$\\
Face $G$ lies in the closure of $F$ & $G\leq F$ in the covector order\\
Hyperplanes containing $G$ & Zero set $z(G)$\\
Codimension of $G$ & $\rk_M z(G)$\\
Nonempty common face of $\bar P,\bar Q$ & Existence of a feasible $G\leq P,Q$\\
Hyperplanes separating $P,Q$ & $\Sep(P,Q)=\{i\mid P(i)\ne Q(i)\}$\\
$\bar P\cap H_J\ne\varnothing$ & Feasible $G\leq P$ with $J\subseteq z(G)$\\
Bounded face $G$ & Every nonzero $F\leq G$ has $F(0)=+$\\
Bounded chamber $P\in\cB$ & Bounded feasible tope\\
Small displacement from $F$ towards $G$ & Composition $F\circ G$\\
\bottomrule
\end{tabularx}
\end{center}
